\chapter{Calculus of Variations and Distributions}

\section{Quadratic minimization}

\begin{notetheorem}[Coercive quadratic minimization]
Let $V$ be a real Banach space and let
$a\colon V\times V\to\mathbb R$ be continuous, symmetric, and coercive:
there are $M,\alpha>0$ such that
\[
  |a(u,v)|\le M\|u\|\,\|v\|,
  \qquad
  a(v,v)\ge\alpha\|v\|^2.
\]
For $\ell\in V^*$ define
\[
  J_\ell(v)=\frac12a(v,v)-\ell(v).
\]
If $K\subseteq V$ is nonempty, closed, and convex, then $J_\ell$ has a unique
minimizer $u\in K$.  The solution map $\ell\mapsto u$ is Lipschitz continuous.
It is linear when $K$ is a linear subspace.
\end{notetheorem}

\begin{proof}
The form $a$ is an inner product, and
$\|v\|_a=\sqrt{a(v,v)}$ is equivalent to the original norm.  Since $V$ is
complete in the original norm, it is complete in $\|\cdot\|_a$ and is
therefore a Hilbert space.  Riesz representation gives a unique $c\in V$ such
that
\[
  \ell(v)=a(c,v).
\]
Then
\[
  J_\ell(v)
  =\frac12\|v-c\|_a^2-\frac12\|c\|_a^2.
\]
Thus the minimizer is the metric projection $P_K^ac$ in the Hilbert norm
$\|\cdot\|_a$.  Existence, uniqueness, and nonexpansiveness follow from the
projection theorem.  To make the dependence on the right-hand side explicit,
let $c_i$ be the Riesz representative of $\ell_i$ with respect to $a$.  Then
\[
  \|c_1-c_2\|_a^2
  =(\ell_1-\ell_2)(c_1-c_2)
  \le \|\ell_1-\ell_2\|_{V^*}\,\|c_1-c_2\|.
\]
Coercivity gives $\|w\|\le\alpha^{-1/2}\|w\|_a$, and hence
\[
  \|c_1-c_2\|_a
  \le \alpha^{-1/2}\|\ell_1-\ell_2\|_{V^*}.
\]
Writing $u_i=P_K^ac_i$ and using nonexpansiveness of the projection once more,
\[
  \|u_1-u_2\|
  \le \alpha^{-1/2}\|u_1-u_2\|_a
  \le \alpha^{-1}\|\ell_1-\ell_2\|_{V^*}.
\]
Thus the solution map is Lipschitz.  If $K$ is a subspace, the orthogonal
projection is linear, so $\ell\mapsto u$ is linear as well.
\end{proof}

\begin{notecorollary}[Variational inequality]
The minimizer $u\in K$ is characterized by
\[
  a(u,v-u)\ge\ell(v-u)
  \qquad(v\in K).
\]
If $K$ is a linear subspace, this is equivalent to the variational equation
\[
  a(u,v)=\ell(v)
  \qquad(v\in K).
\]
\end{notecorollary}

\begin{proof}
For $v\in K$ and $0<t\le1$, minimality of $u$ applied to
$u+t(v-u)$ gives
\[
  0\le\frac{J(u+t(v-u))-J(u)}{t}
  =a(u,v-u)-\ell(v-u)+\frac t2a(v-u,v-u).
\]
Let $t\downarrow0$.  Conversely, expanding $J(v)-J(u)$ and using the
inequality gives nonnegativity.  If $K$ is a subspace, test with
$u\pm v$ to obtain equality.
\end{proof}

\begin{noteremark}
The first variation of
$J(v)=\frac12a(v,v)-\ell(v)$ in the direction $w$ is
\[
  J'(v)w=a(v,w)-\ell(w).
\]
The additional term $\frac12a(w,w)$ in the finite difference
$J(v+w)-J(v)$ is second order.  Quadratic minimization on a subspace therefore
produces a linear variational equation, while nonlinear functionals generally
produce nonlinear Euler--Lagrange equations.
\end{noteremark}

\section{Lax--Milgram and Stampacchia}

\begin{notetheorem}[Lax--Milgram lemma]
Let $H$ be a real Hilbert space and let
$a\colon H\times H\to\mathbb R$ be a continuous bilinear form, not
necessarily symmetric.
Assume coercivity:
\[
  a(v,v)\ge\alpha\|v\|^2
  \qquad(v\in H)
\]
for some $\alpha>0$.  Then, for every $\ell\in H^*$, there is a unique
$u\in H$ such that
\[
  a(u,v)=\ell(v)
  \qquad(v\in H).
\]
Moreover,
\[
  \|u\|\le\frac1\alpha\|\ell\|,
\]
and the map $\ell\mapsto u$ is bounded and linear.
\end{notetheorem}

\begin{proof}
Riesz representation defines a bounded operator $A\in\Bop{H}{H}$ by
\[
  a(u,v)=\inner{Au}{v}.
\]
Coercivity gives
\[
  \alpha\|u\|^2
  \le\inner{Au}{u}
  \le\|Au\|\,\|u\|,
\]
so $A$ is injective, bounded below, and has closed range.  If
$y\perp\Ran(A)$, then $Ay\in\Ran(A)$ and therefore
\[
  \alpha\|y\|^2
  \le a(y,y)
  =\inner{Ay}{y}
  =0.
\]
Thus $y=0$.  Therefore the range is dense and closed, hence all of $H$.  Solve
$Au=f$, where $f$ is the Riesz representative of $\ell$.  The estimate follows
from the lower bound.
\end{proof}

\begin{notetheorem}[Stampacchia variational inequality]
Let $H$ be a real Hilbert space, let $K\subseteq H$ be nonempty, closed, and
convex, and let $a\colon H\times H\to\mathbb R$ be continuous and coercive.
For every
$\ell\in H^*$ there is a unique $u\in K$ such that
\[
  a(u,v-u)\ge\ell(v-u)
  \qquad(v\in K).
\]
The solution depends Lipschitz continuously on $\ell$.
\end{notetheorem}

\begin{proof}
Let $A$ and $f$ be the Riesz representatives of $a$ and $\ell$.  For
$\rho>0$, define
\[
  \Phi(u)=P_K(u-\rho(Au-f)).
\]
For $u,v\in H$, nonexpansiveness of $P_K$ gives
\begin{align*}
  \|\Phi(u)-\Phi(v)\|^2
  &\le
  \|(u-v)-\rho A(u-v)\|^2 \\
  &=\|u-v\|^2
    -2\rho\inner{A(u-v)}{u-v}
    +\rho^2\|A(u-v)\|^2 \\
  &\le
  \bigl(1-2\rho\alpha+\rho^2\|A\|^2\bigr)\|u-v\|^2.
\end{align*}
Choose $0<\rho<2\alpha/\|A\|^2$.  The number
$q^2=1-2\rho\alpha+\rho^2\|A\|^2$ then satisfies $0\le q<1$, so $\Phi$ is a
contraction and has a unique fixed point $u\in K$.

It remains to identify the fixed-point equation.  For $z\in H$, the
characterization of the metric projection says
\[
  u=P_Kz
  \quad\Longleftrightarrow\quad
  \inner{z-u}{v-u}\le0
  \quad(v\in K).
\]
Taking $z=u-\rho(Au-f)$ shows that $u=\Phi(u)$ if and only if
\[
  \inner{Au-f}{v-u}\ge0
  \qquad(v\in K),
\]
which is precisely the stated variational inequality.

Finally, let $u_i$ correspond to $\ell_i$.  Testing the inequality for $u_1$
with $v=u_2$ and that for $u_2$ with $v=u_1$, then adding, gives
\[
  \alpha\|u_1-u_2\|^2
  \le (\ell_1-\ell_2)(u_1-u_2)
  \le\|\ell_1-\ell_2\|\,\|u_1-u_2\|.
\]
After cancelling $\|u_1-u_2\|$ when it is nonzero, one obtains
\[
  \|u_1-u_2\|
  \le\frac1\alpha\|\ell_1-\ell_2\|,
\]
which proves the claimed Lipschitz dependence.
\end{proof}

\section{Test functions and weak derivatives}

\begin{notedefinition}
Let $\Omega\subseteq\mathbb R^N$ be open.  The test-function space is
\[
  \mathcal D(\Omega)=C_c^\infty(\Omega).
\]
The space $L^1_{\mathrm{loc}}(\Omega)$ consists of measurable functions that
are integrable over every compact subset of $\Omega$.
\end{notedefinition}

\begin{notedefinition}
Let $u\in L^1_{\mathrm{loc}}(\Omega)$ and let $\alpha$ be a multi-index.  A
function $v\in L^1_{\mathrm{loc}}(\Omega)$ is the weak derivative
$D^\alpha u$ if
\[
  \int_\Omega v\varphi\,\dd x
  =(-1)^{|\alpha|}
    \int_\Omega uD^\alpha\varphi\,\dd x
  \qquad(\varphi\in\mathcal D(\Omega)).
\]
\end{notedefinition}

\section{Distributions}

\begin{notedefinition}
A \emph{distribution} on $\Omega$ is a linear functional
$T\colon\mathcal D(\Omega)\to\mathbb K$ such that, for every compact
$K\subseteq\Omega$, there are $C_K>0$ and $m_K\in\mathbb N_0$ with
\[
  |T(\varphi)|
  \le C_K
      \max_{|\alpha|\le m_K}
      \sup_{x\in K}|D^\alpha\varphi(x)|
\]
for every $\varphi\in\mathcal D(\Omega)$ supported in $K$.  The space of all
distributions is denoted by $\mathcal D'(\Omega)$.
\end{notedefinition}

\begin{noteremark}[Topology]
A sequence $\varphi_n$ converges to $\varphi$ in $\mathcal D(\Omega)$ if there
is one compact $K\Subset\Omega$ containing all supports and
\[
  \sup_{x\in K}|D^\alpha(\varphi_n-\varphi)(x)|\longrightarrow0
\]
for every multi-index $\alpha$.  The full test-function topology is an
inductive-limit topology and is generally not metrizable.  A sequence
$T_n\in\mathcal D'(\Omega)$ converges distributionally to $T$ if
\[
  T_n(\varphi)\longrightarrow T(\varphi)
  \qquad(\varphi\in\mathcal D(\Omega)).
\]
\end{noteremark}

\begin{noteproposition}[Regular distributions]
Every $u\in L^1_{\mathrm{loc}}(\Omega)$ defines a distribution
$T_u$ by
\[
  T_u(\varphi)=\int_\Omega u\varphi\,\dd x.
\]
The map $u\mapsto T_u$ is injective modulo equality almost everywhere.
\end{noteproposition}

\begin{proof}
If $\supp\varphi\subseteq K$, then
\[
  |T_u(\varphi)|
  \le\|u\|_{L^1(K)}\|\varphi\|_{L^\infty(K)},
\]
so $T_u$ is a distribution of order zero on each compact set.

To prove injectivity, suppose $T_u=0$.  Fix a ball
$B(x_0,2r)\Subset\Omega$ and a standard mollifier $\rho_\varepsilon$ with
$0<\varepsilon<r$.  For every $x\in B(x_0,r)$, the function
\[
  y\longmapsto\rho_\varepsilon(x-y)
\]
belongs to $C_c^\infty(B(x_0,2r))$.  Hence
\[
  (u*\rho_\varepsilon)(x)
  =\int_\Omega u(y)\rho_\varepsilon(x-y)\,\dd y
  =T_u\bigl(\rho_\varepsilon(x-\cdot)\bigr)
  =0.
\]
Local approximation by mollifiers gives
$u*\rho_\varepsilon\to u$ in $L^1(B(x_0,r))$, so $u=0$ almost everywhere on
$B(x_0,r)$.  Since such balls cover $\Omega$, $u=0$ almost everywhere on
$\Omega$.  Thus $u\mapsto T_u$ is injective modulo almost-everywhere equality.
\end{proof}

\begin{noteexample}[Dirac delta]
For $a\in\Omega$, the Dirac distribution is
\[
  \delta_a(\varphi)=\varphi(a).
\]
It is not a regular distribution.
\end{noteexample}

\begin{proof}
Suppose $\delta_a=T_u$ for some $u\in L^1_{\mathrm{loc}}$.  Choose
$\psi\in C_c^\infty(B(0,1))$ with $0\le\psi\le1$ and $\psi(0)=1$, and set
$\psi_k(x)=\psi(k(x-a))$.  Then $\delta_a(\psi_k)=1$, whereas
\[
  \left|\int u\psi_k\right|
  \le\int_{B(a,1/k)}|u(x)|\,\dd x\longrightarrow0
\]
by absolute continuity of the Lebesgue integral, a contradiction.
\end{proof}

\begin{notedefinition}[Distributional derivatives]
For $T\in\mathcal D'(\Omega)$ and a multi-index $\alpha$, define
\[
  D^\alpha T(\varphi)
  =(-1)^{|\alpha|}T(D^\alpha\varphi).
\]
\end{notedefinition}

\begin{proof}[Continuity of the derivative]
Let $K\Subset\Omega$ and suppose that on test functions supported in $K$ the
distribution $T$ satisfies an estimate of order $m_K$.  Since
$\supp(D^\alpha\varphi)\subseteq\supp\varphi$, we have
\[
\begin{aligned}
  |D^\alpha T(\varphi)|
  &=|T(D^\alpha\varphi)| \\
  &\le C_K
      \max_{|\beta|\le m_K}
      \sup_{x\in K}|D^{\alpha+\beta}\varphi(x)| \\
  &\le C_K
      \max_{|\gamma|\le m_K+|\alpha|}
      \sup_{x\in K}|D^\gamma\varphi(x)|.
\end{aligned}
\]
Thus $D^\alpha T$ is a distribution, of local order at most
$m_K+|\alpha|$ on $K$.
\end{proof}

\begin{noteproposition}
If $u\in C^{|\alpha|}(\Omega)$, then
\[
  D^\alpha T_u=T_{D^\alpha u}.
\]
More generally, when $u\in L^1_{\mathrm{loc}}$, the regular distribution
$T_v$ is $D^\alpha T_u$ exactly when $v$ is the weak derivative
$D^\alpha u$.
\end{noteproposition}

\begin{proof}
Repeated integration by parts against compactly supported test functions gives
\[
  D^\alpha T_u(\varphi)
  =(-1)^{|\alpha|}\int uD^\alpha\varphi
  =\int D^\alpha u\,\varphi.
\]
The final statement is precisely the definition of weak differentiation.
\end{proof}
